\documentclass[11pt]{article}
\usepackage[a4paper,margin=26mm]{geometry}
\usepackage{amsmath,amssymb,amsthm,mathtools}
\newtheorem*{theorem}{Theorem}
\newcommand{\Ds}{\mathcal D'}
\newcommand{\id}{\mathbf 1}
\newcommand{\dd}{\partial}
\DeclareMathOperator{\diag}{diag}
\title{All-order continuation of massive free-shell distributions}
\author{}\date{}
\begin{document}
\maketitle
\vspace{-2em}
Let $p_\mu$ be covariant momentum coordinates and $x^\mu$ spacetime coordinates. Use the metric
$\eta^{\mu\nu}=\diag(1,-1,-1,-1)$. Write $p^\mu=\eta^{\mu\nu}p_\nu$,
$\dd_\mu=\dd/\dd x^\mu$, $p^2=p^\mu p_\mu$ and sum repeated indices.
Fix $m\in\mathbb R\setminus\{0\}$. Suppose $\gamma^\mu\gamma^\nu+\gamma^\nu\gamma^\mu
=2\eta^{\mu\nu}\id$, and put $K=\gamma^\mu p_\mu$, $H=K-m\id$.
Take $\mathcal V=M_4(\Ds(\mathbb R^4_x\times\mathbb R^4_p;\mathbb C))$
with the usual smooth compactly supported test functions.
Use the formal Moyal product with first contraction
$\partial_p\otimes\partial_x-\partial_x\otimes\partial_p$ and coefficient $i/2$.

\begin{theorem}
For $F\in\mathcal V$, a sequence $W\in\mathcal V[[\hbar]]$ satisfying
$W|_{\hbar=0}=F$ and $H\star W=W\star H=0$ exists if and only if
\[
 HF=0,\qquad FH=0,\qquad TF=0,\qquad T=p^\mu\dd_\mu .
\]
All three conditions are distribution identities on the leading datum alone.
\end{theorem}

\begin{proof}
Write $L$ and $R$ for left and right star multiplication by $K$. Since $K$ is
linear in momentum and independent of spacetime, every Moyal coefficient of order
at least two acting on either side vanishes. Each monomial of order $r$ differentiates
each factor exactly $r$ times. This gives
\[
 LV=KV+\tfrac{i\hbar}{2}\gamma^\mu\dd_\mu V,\qquad
 RV=VK-\tfrac{i\hbar}{2}(\dd_\mu V)\gamma^\mu .
\]
These operators commute because left and right matrix multiplications commute,
and the spacetime derivatives commute with each other and with momentum
multiplication. With $q=p^2-m^2$ and $B=\eta^{\mu\nu}\dd_\mu\dd_\nu/4$,
the Clifford identity yields
\[
 L^2-m^2=q+i\hbar T-\hbar^2B,\qquad
 R^2-m^2=q-i\hbar T-\hbar^2B .
\]
Indeed the anticommutator of $K$ and $\gamma^\mu$ equals $2p^\mu\id$,
and contraction of two commuting derivatives with $\gamma^\mu\gamma^\nu$
retains only its symmetric part. For any solution the two squared equations
vanish. Their difference yields $2i\hbar TW=0$, hence $TF=0$.
The original equations at order zero give $HF=FH=0$.

Conversely assume the three conditions. From $(K+m)HF=0$ we get $qF=0$.
Let $E=p_\mu\partial_{p_\mu}$, with the derivative applied before multiplication,
and put $D=-E/2$. All these operators act on global distributions. The distributional
product rule gives the following identities, where scalars act by
multiples of the identity.
\begin{equation}
 qD=(D+1)q+m^2,\qquad TD=(D+\tfrac12)T,\qquad
 [D,B]=[B,T]=[B,q]=0 .                                      \tag{1}
\end{equation}
In particular $E(q)=2p^2$ and $[E,T]=T$. These identities involve no
restriction of a distribution to the shell.

Construct $V_0=F$ and, for $j\geq0$, set
\begin{equation}
 V_{j+1}=\frac{(D-j)BV_j}{m^2(j+1)} .                       \tag{2}
\end{equation}
We claim $TV_j=0$ for all $j$ and $qV_{j+1}=BV_j$. Transport follows by
induction from (1). For $j=0$, (1) and $qF=0$ give $qDBF=m^2BF$.
If $j\geq1$ and $qV_j=BV_{j-1}$, then
\begin{align*}
 m^2(j+1)qV_{j+1}
 &= (D-(j-1))BqV_j+m^2 BV_j\\
 &= B(D-(j-1))BV_{j-1}+m^2BV_j
  =m^2(j+1)BV_j,
\end{align*}
where the last equality uses (2) for $V_j$. This proves the claim.

Now take $U=\sum_{j\geq0}\hbar^{2j}V_j$. The claim gives $TU=0$ and
$(q-\hbar^2 B)U=0$. Thus $(L^2-m^2)U=(R^2-m^2)U=0$. Define
\[
 W=\frac{(L+m)(R+m)}{4m^2}\,U .
\]
Commutation of $L$ and $R$ gives $(L-m)W=(R-m)W=0$. At order
zero, $(K+m)F(K+m)/(4m^2)=F$ follows from $KF=FK=mF$.
Each coefficient uses finitely many derivatives and smooth polynomial
multipliers of $F$, so it is a distribution on the original full phase space.
\end{proof}
\paragraph{A spatially varying example.}
At $p_*=(m,0,0,0)$ let $P=(\id+\gamma^0)/2$. The leading datum
$F=\cos(x^1)P\,\delta(p-p_*)$ obeys both ordinary shell equations. Its
transport equation reduces to $m\partial_0F=0$. Since $U$ has no odd powers,
the first coefficient of the constructed solution is
\[
 W_1=\frac{i}{4m}[\gamma^\mu,\partial_\mu F]
    =-\frac{i}{4m}\sin(x^1)\,\gamma^1\gamma^0\,\delta(p-p_*).
\]
It is nonzero since $(\gamma^1\gamma^0)^2=\id$. Thus even smooth variation of
the internal density along a spatial coordinate can require a correction.

\noindent
The result permits spacetime dependence and distributions on both massive sheets.
The corrections can contain normal derivatives of shell distributions. The
construction gives a linear choice of continuation. For $F\ne0$, multiplication
of a solution by $1+c\hbar$ already gives distinct continuations with the same
leading datum as $c$ varies. Existence here concerns formal series for the free
generator.
\end{document}
